% ============================================================
% CAPÍTULO 2: ARREGLO EXPERIMENTAL
% ============================================================
\chapter{Arreglo experimental}
\label{cap:arreglo}

En este capítulo se describe el arreglo experimental empleado para la
medición fotoacústica de líquidos. Se presentan la celda que contiene
la muestra, los equipos de medición involucrados, la fuente de
excitación y su trayectoria óptica, y los tres subsistemas de
instrumentación que fueron automatizados en este trabajo: el
osciloscopio, el láser y el módulo de medición de temperatura.

El arreglo parte de una configuración experimental preexistente en el
Laboratorio de Fotofísica y Películas Delgadas del ICAT. La aportación
de este trabajo no consiste en modificar dicha configuración, sino en
dotarla de control programático y adquisición sincronizada, de modo
que las secciones \ref{sec:celda} a \ref{sec:excitacion} describen el
arreglo heredado y las secciones \ref{sec:auto_osci} a
\ref{sec:modulo_temperatura} describen los subsistemas desarrollados.

% ============================================================
\section{Descripción de la celda fotoacústica}
\label{sec:celda}
% ============================================================

La celda que contiene la muestra líquida fue diseñada y manufacturada
por Rangel Mendieta como parte del arreglo experimental desarrollado
para la caracterización del hidrófono empleado en este trabajo
\cite{Rangel2023}. En la presente investigación se reutiliza la misma
celda física, por lo que su geometría constituye una condición de
frontera del experimento y no un parámetro de diseño.

\subsection*{Construcción y dimensiones}

La celda es un recipiente prismático de acrílico de \SI{3}{\milli\meter}
de espesor, ensamblado mediante pestañas y ranuras y sellado con
adhesivo. Incorpora dos ventanas de cuarzo que permiten el paso del haz
láser sin las pérdidas ni la fluorescencia que introduciría el
acrílico. La \cref{tab:piezas_celda} resume las piezas que la componen.

\begin{table}[H]
  \centering
  \caption{Piezas que integran la celda fotoacústica \cite{Rangel2023}.}
  \label{tab:piezas_celda}
  \begin{tabular}{clcc}
    \toprule
    \textbf{Pieza} & \textbf{Función} & \textbf{Dimensiones [\si{\milli\meter}]}
    & \textbf{Material} \\
    \midrule
    1 & Base con ranuras de ensamble & $220 \times 150 \times 3$ & Acrílico \\
    2 & Paredes largas (dos piezas) & $206 \times 70 \times 3$ & Acrílico \\
    3 & Placa de sujeción, barreno \diameter\,12 (cabezal) & $76 \times 76 \times 3$ & Acrílico \\
    4 & Placa de sujeción, barreno \diameter\,16 (cuerpo) & $76 \times 76 \times 3$ & Acrílico \\
    5 & Placa de sujeción, barreno \diameter\,24 (cable) & $76 \times 76 \times 3$ & Acrílico \\
    6 & Ventana rectangular & $54 \times 16 \times 0.5$ & Cuarzo \\
    7 & Ventana cuadrada & $18 \times 18 \times 0.5$ & Cuarzo \\
    \bottomrule
  \end{tabular}
\end{table}

La tabla reproduce las piezas reportadas por el autor de la celda y no
constituye una lista exhaustiva de materiales: la pared transversal
opuesta a la del hidrófono no aparece detallada por separado en esa
fuente. De las dos ventanas de cuarzo, la rectangular es la que aloja
el ingreso del haz en una de las paredes largas, según se describe en
la \cref{sec:excitacion}.

El ensamble resulta en una cavidad interna aproximada de
$200 \times 70 \times \SI{70}{\milli\meter}$. El volumen de muestra
empleado es de \SI{400}{\milli\liter} \cite{Rangel2023}, que
corresponde a un tirante de líquido de
\begin{equation}
  h = \frac{V}{A_\text{base}}
    = \frac{\SI{400}{\centi\meter\cubed}}
           {\SI{20.0}{\centi\meter} \times \SI{7.0}{\centi\meter}}
    \approx \SI{28.6}{\milli\meter}.
  \label{eq:tirante}
\end{equation}

Este valor no es arbitrario. La ranura que aloja la ventana
rectangular de cuarzo en la pared larga abarca de \SI{12}{\milli\meter}
a \SI{29}{\milli\meter} de altura medidos desde el fondo de la celda,
de modo que el volumen de \SI{400}{\milli\liter} corresponde
precisamente al nivel mínimo que cubre por completo la ventana de
entrada del haz. Un volumen menor dejaría parte de la ventana en
contacto con aire y modificaría la trayectoria óptica; uno mayor no
aporta beneficio y aproxima la superficie libre al eje del haz.

% \begin{figure}[H]
%   \centering
%   \includegraphics[width=0.8\textwidth]{celda_ensamble}
%   \caption{Vista isométrica de la celda fotoacústica ensamblada.}
%   \label{fig:celda_ensamble}
% \end{figure}

\subsection*{Montaje del hidrófono}

Una de las dos paredes transversales de la celda está formada por tres
placas de acrílico dispuestas coaxialmente, con barrenos de
\diameter\,12, \diameter\,16 y \diameter\,24\,\si{\milli\meter}
que alojan sucesivamente el cabezal, el cuerpo y el cable del
transductor. El hidrófono empleado es el desarrollado por Rangel
Mendieta \cite{Rangel2023}, cuyas características se detallan en la
\cref{sec:equipos}.

El escalonamiento de los tres diámetros no obedece únicamente a la
sujeción: al apoyar el transductor en tres secciones separadas a lo
largo de su eje, el conjunto restringe su inclinación y mantiene la
cara sensora perpendicular al eje del haz. Esa perpendicularidad es la
condición que permite tratar la separación entre el hidrófono y el haz
como una distancia única, según se detalla en el
\cref{ap:tiempos_arribo}. El montaje garantiza además que la posición
relativa entre el hidrófono y la pared de la celda se conserve entre
sesiones experimentales.

\subsection*{Geometría de la interacción y consecuencias acústicas}

El haz láser ingresa a la muestra a través de la ventana rectangular
de cuarzo situada en una de las paredes largas, de manera que su
trayectoria dentro del líquido corresponde al ancho interno de la
cavidad, es decir, aproximadamente \SI{70}{\milli\meter}. El haz
viaja paralelo a la cara sensora del hidrófono, a una distancia
aproximada de \SI{10}{\milli\meter} \cite{Rangel2023}.

Esta configuración reproduce la geometría de fuente cilíndrica
descrita en la \cref{sec:fotoacustico_liquidos}: una columna
iluminada larga y delgada que emite una onda cilíndrica hacia un
transductor situado a corta distancia perpendicular
\cite{PatelTam1981}.

Ahora bien, las dimensiones finitas de la celda imponen un límite
temporal a la señal utilizable. Tomando la velocidad del sonido en
agua como \SI{1480}{\meter\per\second} a \SI{20}{\celsius}, los
tiempos de arribo al hidrófono de la señal directa y de las primeras
reflexiones en las fronteras de la cavidad son los indicados en la
\cref{tab:tiempos_celda}. El desarrollo del cálculo, junto con las
hipótesis geométricas empleadas, se presenta en el
\cref{ap:tiempos_arribo}.

\begin{table}[H]
  \centering
  \caption{Tiempos de arribo estimados al hidrófono para la señal
           directa y las primeras reflexiones en las fronteras de la
           celda.}
  \label{tab:tiempos_celda}
  \begin{tabular}{lcc}
    \toprule
    \textbf{Trayectoria} & \textbf{Distancia [\si{\milli\meter}]}
    & \textbf{Arribo [\si{\micro\second}]} \\
    \midrule
    Señal directa                      &  10 &  6.8 \\
    Reflexión en superficie libre      &  19 & 12.9 \\
    Reflexión en el fondo              &  42 & 28.5 \\
    Reflexión en paredes transversales & 100 & 67.6 \\
    \bottomrule
  \end{tabular}
\end{table}

De la \cref{tab:tiempos_celda} se desprende que la ventana temporal
libre de reflexiones es de aproximadamente \SI{6}{\micro\second}
contados a partir de la llegada de la señal directa. Este intervalo
es el que contiene el pulso bipolar de generación fotoacústica en
volumen, y determina tanto la base de tiempo que debe configurarse
en el osciloscopio como la longitud de registro necesaria para
resolverlo adecuadamente. Quan \textit{et al.} advierten
explícitamente que las dimensiones finitas del recipiente producen
reflexiones acústicas múltiples que se superponen a la señal de
interés y complican su interpretación \cite{Quan1994}; la
\cref{tab:tiempos_celda} cuantifica ese efecto para la celda
particular empleada en este trabajo.

\subsection*{Naturaleza de los tiempos de la tabla \ref{tab:tiempos_celda}}
\label{sec:origen_tiempos}

Los valores de la \cref{tab:tiempos_celda} son tiempos de vuelo:
miden el trayecto acústico desde la columna iluminada hasta el
transductor, y su origen es el instante de emisión del pulso láser. No
deben confundirse con las posiciones que esos mismos eventos ocupan
dentro de un registro adquirido, y la distinción importa porque el
\cref{cap:resultados} reporta instantes que difieren de ellos en varias
decenas de microsegundos.

La razón es instrumental. El eje temporal reconstruido mediante la
\cref{eq:conv_tiempo} tiene su cero en el evento de disparo, y la
posición de ese evento dentro del registro la fija el operador con el
control de posición horizontal del osciloscopio al preparar la sesión.
Ese control no se opera desde la aplicación, según la decisión de
diseño expuesta en la \cref{sec:decisiones}, y su ajuste no queda
registrado como un desplazamiento conocido. En consecuencia, el
instante absoluto en que aparece la señal dentro de un registro
depende del montaje del día y no es comparable entre sesiones.

Lo que sí es comparable, y lo que la \cref{tab:tiempos_celda} predice,
son las \emph{diferencias} entre arribos: el intervalo de
aproximadamente \SI{6}{\micro\second} que separa la señal directa de
la primera reflexión se conserva con independencia de dónde caiga el
conjunto dentro del registro. Por esa razón, las magnitudes que el
\cref{cap:resultados} emplea para juzgar la repetibilidad del sistema
son intervalos y dispersiones, y no posiciones absolutas sobre el eje
de tiempos.

\subsection*{Condiciones de operación y limitaciones identificadas}

Debe señalarse que la primera reflexión que limita la ventana útil es
la de la superficie libre, cuya posición depende del volumen vertido
según la \cref{eq:tirante}. Una variación en el llenado desplaza este
tiempo de arribo, por lo que el volumen de muestra constituye un
parámetro experimental que debe controlarse y registrarse.

El láser de excitación es un recurso compartido entre varios usuarios
del laboratorio. En consecuencia, ni la celda ni el cabezal del láser
se encuentran fijados de manera permanente a la mesa óptica, y la
geometría del arreglo debe restablecerse al inicio de cada sesión
experimental. Esta condición no afecta la validez de las mediciones
realizadas dentro de una misma sesión, durante la cual el arreglo
permanece inalterado, pero sí impide comparar amplitudes absolutas
entre sesiones distintas sin información adicional.

La estrategia adoptada consiste en registrar la geometría efectiva
como metadato de cada sesión: altura del eje del haz sobre el fondo
interno, distancia entre el hidrófono y el eje del haz, posición del
hidrófono a lo largo de la celda, volumen vertido y temperatura
ambiente inicial. Con estos parámetros, los tiempos de arribo de la
\cref{tab:tiempos_celda} pueden recalcularse para cada sesión
mediante el procedimiento del \cref{ap:tiempos_arribo}, y cualquier
discrepancia entre sesiones puede atribuirse a su causa geométrica.

El uso de la celda a lo largo del desarrollo experimental permitió
identificar además dos limitaciones de su diseño actual.

La primera es el sellado adhesivo entre piezas, que constituye un
punto de falla ante el llenado y vaciado repetidos que exige el
protocolo de medición, y que introduce un material adicional en
contacto con la muestra.

La segunda es el tirante de líquido reducido. Un tirante de
\SI{28.6}{\milli\meter} aproxima la superficie libre al eje del haz y
acorta la ventana temporal disponible; además, favorece la
estratificación térmica vertical del líquido en reposo, aspecto
relevante para la medición de temperatura descrita en la
\cref{sec:modulo_temperatura}.

Ambas observaciones, junto con la conveniencia de fijar la geometría
del arreglo, motivan las propuestas de trabajo a futuro presentadas
en el \cref{cap:conclusiones}.

% ============================================================
\section{Equipos de medición: osciloscopio, hidrófono y fotodiodo}
\label{sec:equipos}
% ============================================================

\subsection*{Osciloscopio}

La adquisición de señales se realiza con un osciloscopio digital de
la serie Tektronix TDS5000B. A lo largo del desarrollo experimental
se emplearon indistintamente los modelos TDS5052B y TDS5054B, cuya
única diferencia relevante es el número de canales de entrada: dos en
el primero y cuatro en el segundo. Dado que el arreglo utiliza
únicamente dos canales --- uno para el hidrófono y otro para el
fotodiodo de disparo --- ambos modelos son equivalentes para los
fines de este trabajo, y el software desarrollado opera sobre
cualquiera de los dos sin modificación.

Los instrumentos de esta serie ejecutan internamente un sistema
operativo Windows y exponen una interfaz de red que implementa el
protocolo VXI-11, lo que permite su control remoto mediante comandos
SCPI sobre Ethernet. Las características de esta interfaz y su
aprovechamiento se describen en la \cref{sec:auto_osci}.

\subsection*{Hidrófono}

El transductor acústico es el hidrófono diseñado, manufacturado y
caracterizado por Rangel Mendieta \cite{Rangel2023}. Su elemento
sensor es una cerámica piezoeléctrica de composición KNLNS-BaCaZr de
\diameter\,10.71\,\si{\milli\meter} y \SI{1.06}{\milli\meter} de
espesor, alojada en una carcasa de uretano de
\diameter\,23\,\si{\milli\meter} exterior. El uretano fue
seleccionado por su impedancia acústica próxima a la del agua, lo que
reduce las pérdidas por desacoplamiento en la interfaz entre el medio
y el transductor.

La caracterización reportada por su autor arroja los resultados
resumidos en la \cref{tab:hidrofono}.

\begin{table}[H]
  \centering
  \caption{Características del hidrófono empleado, según la
           caracterización de su autor \cite{Rangel2023}.}
  \label{tab:hidrofono}
  \begin{tabular}{ll}
    \toprule
    \textbf{Parámetro} & \textbf{Valor} \\
    \midrule
    Material del elemento sensor & Cerámica KNLNS-BaCaZr \\
    Dimensiones del elemento sensor & \diameter\,10.71 $\times$ \SI{1.06}{\milli\meter} \\
    Material de la carcasa & Uretano \\
    Respuesta en frecuencia de la cerámica aislada
      & \SIrange{320}{1582}{\kilo\hertz} \\
    Respuesta en frecuencia del hidrófono ensamblado
      & \SIrange{41}{797.5}{\kilo\hertz} \\
    Frecuencia dominante en agua & \SI{41}{\kilo\hertz} \\
    Resonancias secundarias en agua & \SI{323}{\kilo\hertz}, \SI{797.5}{\kilo\hertz} \\
    Amplitud máxima en agua tridestilada & \SI{50}{\milli\volt} \\
    \bottomrule
  \end{tabular}
\end{table}

Rangel Mendieta reporta que la amplitud de la respuesta fotoacústica
del hidrófono sumergido en agua supera al doble de la obtenida con el
hidrófono en aire, lo que confirma un acoplamiento acústico adecuado
con el medio en el que opera \cite{Rangel2023}.

\subsection*{Consideración sobre la respuesta en frecuencia}

Los valores de la \cref{tab:hidrofono} imponen una restricción que
debe hacerse explícita, pues condiciona la interpretación de las
señales adquiridas.

La ventana temporal libre de reflexiones establecida en la
\cref{sec:celda} es de aproximadamente \SI{6}{\micro\second}. Un
transitorio contenido en ese intervalo posee contenido espectral
significativo por encima de \SI{170}{\kilo\hertz}, y el pulso de
excitación empleado, de \SI{7}{\nano\second} de duración, es aún más
rápido. Sin embargo, la respuesta dominante del hidrófono ensamblado
se localiza en \SI{41}{\kilo\hertz}, cuyo periodo de
\SI{24}{\micro\second} excede la ventana útil completa.

Se concluye que la señal registrada no corresponde a la forma
temporal del pulso fotoacústico, sino predominantemente a la
respuesta propia del transductor excitada por dicho pulso. Es
pertinente notar que la cerámica aislada exhibe una respuesta mucho
más rápida (\SIrange{320}{1582}{\kilo\hertz}) que el conjunto
ensamblado, lo que sugiere que el encapsulado de uretano es el
elemento que limita el ancho de banda del sistema.

Esta observación coincide con la advertencia de Quan \textit{et al.},
quienes señalan que el tiempo de subida del detector queda
determinado por el espesor del elemento piezoeléctrico y que un
detector cuyas dimensiones excedan la longitud de onda del pulso
degrada tanto la resolución espacial como la temporal de la forma de
onda \cite{Quan1994}.

La consecuencia metodológica es que el análisis espectral de las
señales adquiridas caracteriza conjuntamente al líquido y al
transductor, y que la comparación entre líquidos distintos solo es
válida en tanto el transductor permanezca invariante. El desarrollo
de un hidrófono de mayor ancho de banda constituye una línea de
trabajo identificada y se aborda en el \cref{cap:conclusiones}.

\subsection*{Fotodiodo}

El disparo de la adquisición se obtiene mediante un fotodiodo
amplificado de silicio Thorlabs PDA10A \cite{Rangel2023}, que capta
una fracción del pulso láser y entrega un pulso eléctrico al segundo
canal del osciloscopio.

Su función en el arreglo es doble. En primer lugar, provee la señal
de disparo (\textit{trigger}) que sincroniza la adquisición con la
emisión del pulso láser. En segundo lugar, establece el origen
temporal $t_0$ respecto del cual se miden los tiempos de vuelo
acústicos; sin esta referencia, los valores de la
\cref{tab:tiempos_celda} no serían verificables experimentalmente,
ya que no existiría un instante conocido desde el cual contar.

Debe precisarse que el fotodiodo no se emplea como medidor de energía
del pulso. El arreglo no incorpora un sensor calibrado de energía,
por lo que las señales adquiridas no se normalizan respecto de $E_0$.
De acuerdo con la \cref{eq:presion_general}, la amplitud registrada
es proporcional al producto $E_0\,\alpha$, de modo que las
variaciones de amplitud entre mediciones no pueden atribuirse de
manera unívoca al coeficiente de absorción del medio. Por esta razón,
el análisis de las señales en este trabajo se sustenta en su forma
temporal y en su contenido espectral, y no en su amplitud absoluta.

% ============================================================
\section{Fuente de excitación y trayectoria óptica}
\label{sec:excitacion}
% ============================================================

\subsection*{Láser de excitación}

La excitación se realiza con un láser pulsado de estado sólido
EKSPLA NL303HT-10-SH, de Nd:YAG con conmutación Q, cuyas
características relevantes se resumen en la \cref{tab:laser}.

\begin{table}[H]
  \centering
  \caption{Características del láser de excitación.}
  \label{tab:laser}
  \begin{tabular}{ll}
    \toprule
    \textbf{Parámetro} & \textbf{Valor} \\
    \midrule
    Modelo & EKSPLA NL303HT-10-SH \\
    Medio activo & Nd:YAG \\
    Longitudes de onda & \SI{1064}{\nano\meter} (fundamental),
                         \SI{532}{\nano\meter} (segundo armónico) \\
    Longitud de onda empleada & \SI{1064}{\nano\meter} \\
    Duración de pulso & \SI{7}{\nano\second} \\
    Frecuencia de repetición & \SI{10}{\hertz} \\
    Interfaz de control & RS-232 \\
    \bottomrule
  \end{tabular}
\end{table}

Las mediciones de este trabajo se realizaron con la línea fundamental
de \SI{1064}{\nano\meter}. La frecuencia de repetición de
\SI{10}{\hertz} es fija y constituye el parámetro sobre el que se
sustenta la estimación del número de pulsos integrados descrita en el
\cref{cap:resultados}.

\subsection*{Control de energía}

El láser no expone un valor de energía por pulso en unidades físicas,
sino un conjunto de niveles discretos de salida denominados
\textit{Output level}, con los valores \textit{OFF},
\textit{Adjustment} y \textit{Max}, complementados por el retardo del
conmutador electro-óptico (\textit{EO delay}), que admite ajuste
numérico. La combinación de ambos parámetros determina la energía
efectivamente emitida.

En consecuencia, la interfaz desarrollada permite seleccionar y
modificar el nivel de salida del láser, pero no cuantificar la
energía resultante. Esta distinción es relevante y se retoma en la
discusión de resultados del \cref{cap:resultados}.

\subsection*{Trayectoria óptica}

El haz emitido incide sobre la ventana rectangular de cuarzo de la
celda, situada a aproximadamente \SI{27.5}{\centi\meter} de la salida
del láser \cite{Rangel2023}, e ingresa a la muestra líquida
atravesando el ancho interno de la cavidad. Una fracción del haz se
deriva hacia el fotodiodo para generar la señal de disparo.

Como se indicó en la \cref{sec:celda}, la trayectoria óptica no se
encuentra fijada de manera permanente y debe restablecerse al inicio
de cada sesión.

% \begin{figure}[H]
%   \centering
%   \includegraphics[width=0.9\textwidth]{esquema_optico}
%   \caption{Esquema del arreglo óptico: láser, fotodiodo, celda
%            fotoacústica y osciloscopio.}
%   \label{fig:esquema_optico}
% \end{figure}

% \begin{figure}[H]
%   \centering
%   \includegraphics[width=0.9\textwidth]{arreglo_fotografia}
%   \caption{Fotografía del arreglo experimental montado.}
%   \label{fig:arreglo_fotografia}
% \end{figure}

% ============================================================
\section{Automatización del osciloscopio}
\label{sec:auto_osci}
% ============================================================

\subsection*{Interfaz de comunicación}

Los osciloscopios de la serie TDS5000B implementan el protocolo
VXI-11, un estándar de control de instrumentos sobre TCP/IP basado en
llamadas a procedimiento remoto. La comunicación se establece con la
biblioteca \texttt{python-vxi11}, que expone las operaciones de
escritura y consulta necesarias para enviar comandos SCPI y recibir
sus respuestas.

Se descartó el uso de PyVISA, alternativa habitual para el control de
instrumentos, debido a que su capa de abstracción requiere una
implementación de VISA instalada en el sistema anfitrión y no aportó
ventaja alguna frente al acceso directo al protocolo. La biblioteca
\texttt{python-vxi11} implementa VXI-11 de forma nativa en Python y
elimina esa dependencia externa.

Una particularidad de estos instrumentos condiciona su operación
remota: dado que ejecutan internamente un sistema operativo Windows,
el servidor VXI-11 solo responde mientras la aplicación nativa del
osciloscopio se encuentra en ejecución. Si dicha aplicación está
cerrada, los intentos de conexión fallan aunque el instrumento
responda a nivel de red.

\subsection*{Configuración de la transferencia de datos}

Antes de cualquier adquisición, el módulo de control establece los
parámetros que gobiernan el formato de transferencia de la forma de
onda, mediante la secuencia de comandos SCPI de la
\cref{lst:config_base}.

\begin{lstlisting}[caption={Configuración base de la transferencia de
    forma de onda.},
    label={lst:config_base},
    language={}]
DATA:SOURCE CH1
DATA:ENCDG RIBINARY
DATA:WIDTH 2
DATA:START 1
DATA:STOP 1000000
WFMPRE:NR_PT?
\end{lstlisting}

La codificación \texttt{RIBINARY} con ancho de dos bytes transfiere
las muestras como enteros con signo en formato binario de orden de
byte grande, lo que reduce el volumen de datos y el tiempo de
transferencia respecto de una codificación en texto.

El valor de \texttt{DATA:STOP} se establece deliberadamente por
encima de cualquier longitud de registro admitida por el instrumento;
el osciloscopio lo limita internamente a la longitud real
configurada. La consulta posterior de \texttt{WFMPRE:NR\_PT?} devuelve
el número efectivo de puntos que serán transferidos. Este orden de
operaciones es indispensable: fijar \texttt{DATA:STOP} a un valor
inferior al desplazamiento del punto de disparo provoca que la
transferencia contenga únicamente muestras anteriores al evento de
interés, sin que el instrumento reporte error alguno.

\subsection*{Conversión a unidades físicas}

La respuesta al comando \texttt{CURVE?} consiste en un bloque binario
precedido por una cabecera que indica su longitud, y contiene los
valores digitalizados sin escalar. Su conversión a unidades físicas
requiere los factores de escala que el instrumento expone a través
del preámbulo \texttt{WFMPRE}: \texttt{YMULT}, \texttt{YOFF},
\texttt{YZERO}, \texttt{XINCR}, \texttt{XZERO} y \texttt{PT\_OFF}.

El voltaje correspondiente a la muestra $i$-ésima se obtiene como
\begin{equation}
  V_i = \left(D_i - \texttt{YOFF}\right)\cdot \texttt{YMULT}
        + \texttt{YZERO},
  \label{eq:conv_voltaje}
\end{equation}
donde $D_i$ es el valor entero transferido. El instante asociado a esa
misma muestra es
\begin{equation}
  t_i = \left(i - \texttt{PT\_OFF}\right)\cdot \texttt{XINCR}
        + \texttt{XZERO},
  \label{eq:conv_tiempo}
\end{equation}
en la que \texttt{PT\_OFF} representa el índice de la muestra que
coincide con el evento de disparo, de modo que los instantes
anteriores al mismo resultan negativos.

Ambas conversiones se aplican de forma vectorizada sobre el arreglo
completo mediante \texttt{numpy}, lo que evita iterar sobre registros
que pueden alcanzar varios millones de muestras.

\subsection*{Parámetros de adquisición: lectura y escritura}

El módulo distingue dos grupos de parámetros según quién los
establece.

El primer grupo describe cómo el instrumento digitaliza la señal: la
longitud de registro (\texttt{HOR:RECLENGTH}), la escala vertical
(\texttt{CH$n$:SCALE}), la escala horizontal
(\texttt{HORizontal:SCAle}), el acoplamiento de entrada
(\texttt{CH$n$:COUPling}) y el nivel de disparo
(\texttt{TRIGger:MAIn:LEVel}). Estos parámetros se configuran
físicamente en el panel del osciloscopio y el módulo únicamente los
consulta, sin escribirlos en ningún momento. Los valores leídos se
muestran en la interfaz en modo de solo lectura y se guardan junto con
cada captura. La justificación de esta decisión, que es de
trazabilidad y no de comodidad, se desarrolla en la
\cref{sec:decisiones}.

El segundo grupo gobierna el modo de acumulación y sí se escribe desde
la aplicación: el modo de adquisición (\texttt{ACQ:MODE}), el número
de promedios (\texttt{ACQ:NUMAVG}) y las órdenes de inicio y detención
descritas a continuación. Son los parámetros que la secuencia
automática necesita manipular para operar, y ninguno de ellos altera
la calibración de los ejes con la que se reconstruye la señal.

\subsection*{Modos de adquisición}

El módulo implementa dos formas de adquisición que corresponden a los
modos de operación descritos en el \cref{cap:resultados}.

En el modo de adquisición por secuencia se configura
\texttt{ACQ:MODE AVERAGE} junto con
\texttt{ACQ:STOPAFTER SEQUENCE}, de manera que el instrumento
promedia el número de disparos indicado por \texttt{ACQ:NUMAVG} y
detiene la adquisición automáticamente al completarlos. La aplicación
consulta periódicamente \texttt{ACQ:STATE?} hasta constatar su
terminación.

En el modo de adquisición controlada por evento externo, la
aplicación inicia y detiene la integración de forma explícita
mediante \texttt{ACQ:STATE RUN} y \texttt{ACQ:STATE STOP}, en función
de una condición evaluada por el software. El osciloscopio actúa
entonces como integrador y el número de disparos promediados queda
determinado por la duración del intervalo, no por un valor
preestablecido.

\subsection*{Robustez de la comunicación}

Toda operación de lectura se reintenta un número acotado de veces
antes de declararse fallida. El módulo distingue dos categorías de
fallo con tratamiento diferenciado: los errores de conexión, que
invalidan la sesión VXI-11 y marcan el instrumento como desconectado,
y los errores de captura, en los que la conexión permanece activa
pero una transferencia concreta no pudo completarse. Esta distinción
evita que una lectura fallida aislada provoque la desconexión
innecesaria del instrumento y la interrupción de una secuencia en
curso.

Al concluir cualquier operación de adquisición, el módulo restablece
el instrumento a su estado de barrido continuo mediante
\texttt{ACQ:STOPAFTER RUNSTOP} seguido de \texttt{ACQ:STATE RUN}, de
modo que el osciloscopio quede en condiciones de uso manual sin
intervención adicional.

% ============================================================
\section{Automatización del láser}
\label{sec:auto_laser}
% ============================================================

\subsection*{Interfaz de comunicación}

El láser expone su control remoto a través de un enlace RS-232 y de
una biblioteca de enlace dinámico proporcionada por el fabricante,
\texttt{REMOTECONTROL64.dll}, acompañada de un archivo de
configuración que declara los registros accesibles del instrumento y
sus valores admisibles.

La biblioteca se carga desde Python mediante \texttt{ctypes.WinDLL}.
Al tratarse de una biblioteca exclusivamente de \SI{64}{bits}, debe
cargarse en el espacio de proceso de un intérprete de la misma
arquitectura; el uso de Python de \SI{64}{bits} satisface esta
condición sin requerir un proceso puente ni capa intermedia alguna.

La carga presenta una particularidad: la biblioteca resuelve sus
dependencias y localiza su archivo de configuración por rutas
relativas al directorio de trabajo. Por ello, el módulo registra
temporalmente el directorio de la biblioteca como ruta de búsqueda y
cambia a él durante la carga, restaurando el directorio original
inmediatamente después.

\subsection*{Modelo de registros}

El control del instrumento se ejerce mediante lectura y escritura de
registros identificados por nombre, a través de las funciones
\texttt{rcSetRegFromString} y \texttt{rcGetRegAsString}. Los
registros empleados en este trabajo se listan en la
\cref{tab:registros_laser}.

\begin{table}[H]
  \centering
  \caption{Registros del láser empleados por el sistema de control.}
  \label{tab:registros_laser}
  \begin{tabularx}{\textwidth}{lXl}
    \toprule
    \textbf{Registro} & \textbf{Función} & \textbf{Acceso} \\
    \midrule
    \texttt{State} & Inicio y detención de la emisión & Lectura/escritura \\
    \texttt{Output level} & Nivel de energía de salida & Lectura/escritura \\
    \texttt{Continuous / Burst mode / Trigger burst}
      & Modo de disparo & Lectura/escritura \\
    \texttt{Burst length, pulses} & Pulsos por ráfaga & Lectura/escritura \\
    \texttt{Adjustment EO delay} & Retardo del conmutador electro-óptico
      & Lectura/escritura \\
    \texttt{Set cooling temperature}& Temperatura objetivo del refrigerante
      & Lectura/escritura \\
    \texttt{Read cooling temperature}
      & Temperatura actual del refrigerante
      & Solo lectura \\
    \texttt{Lamp pulse counter} & Contador acumulado de disparos & Solo lectura \\
    \bottomrule
  \end{tabularx}
\end{table}

El módulo encapsula el acceso a estos registros tras una interfaz de
métodos con nombres descriptivos, de modo que ni la interfaz gráfica
ni el orquestador de mediciones manipulan cadenas de registro
directamente. Todas las llamadas a la biblioteca se serializan
mediante un mecanismo de exclusión mutua, requisito impuesto por el
carácter no reentrante del enlace RS-232 subyacente.

\subsection*{Procedimiento de estado seguro}

El sistema incorpora un procedimiento que devuelve el láser a una
condición de operación segura. Dicho procedimiento se activa ante
cuatro condiciones: conclusión normal de una secuencia de medición,
parada de emergencia solicitada por el operador, fallo irrecuperable
de comunicación tras agotar los reintentos de reconexión, e
inactividad prolongada en modo manual.

El procedimiento ejecuta en orden los cuatro comandos de la
\cref{tab:modo_seguro}.

\begin{table}[H]
  \centering
  \caption{Secuencia de comandos del procedimiento de estado seguro.}
  \label{tab:modo_seguro}
  \begin{tabular}{clc}
    \toprule
    \textbf{Orden} & \textbf{Acción} & \textbf{Registro} \\
    \midrule
    1 & Detener la emisión & \texttt{State $\rightarrow$ STOP} \\
    2 & Establecer el retardo EO en 3800 & \texttt{Adjustment EO delay} \\
    3 & Anular el nivel de salida & \texttt{Output level $\rightarrow$ OFF} \\
    4 & Restituir el modo de disparo continuo
      & \texttt{Continuous / Burst mode} \\
    \bottomrule
  \end{tabular}
\end{table}

El procedimiento intenta la totalidad de los comandos con
independencia del resultado de cada uno, en lugar de abortar ante el
primer fallo. Esta decisión de diseño obedece a que un fallo de
comunicación parcial no debe impedir que se ejecuten los comandos
restantes: el objetivo es maximizar el número de acciones de
seguridad efectivamente aplicadas. El resultado individual de cada
comando se reporta a la interfaz, que advierte al operador si alguno
no pudo confirmarse.

Debe subrayarse que el procedimiento actúa exclusivamente sobre el
láser. Los parámetros del osciloscopio no se modifican en ningún
caso, de modo que la configuración de adquisición establecida por el
operador se conserva íntegra tras un evento de seguridad.

% ============================================================
\section{Módulo de medición de temperatura}
\label{sec:modulo_temperatura}
% ============================================================

La medición de la temperatura de la muestra constituye un subsistema
desarrollado íntegramente en este trabajo. Su necesidad se estableció
en la \cref{sec:fotoacustico_liquidos}: la señal fotoacústica depende
de $\beta$, $C_p$ y $v_s$, magnitudes que varían con la temperatura,
y el peso relativo de los mecanismos de generación también lo hace
\cite{Astrath2023}. Registrar la señal sin registrar simultáneamente
la temperatura del medio impide interpretar correctamente las
variaciones observadas.

\subsection*{Arquitectura del módulo}

El módulo se compone de un microcontrolador ESP32 WROOM-32 y cuatro
sensores digitales de temperatura DS18B20 en encapsulado sumergible.
Cada sensor se conecta a un bus 1-Wire independiente, sobre las líneas
de entrada y salida de propósito general 16, 17, 18 y 19 del
microcontrolador, con alimentación externa explícita y una resistencia
de polarización de \SI{4.7}{\kilo\ohm} por bus. La alternativa
habitual ---un único bus compartido por los cuatro sensores--- reduce
el cableado, pero hace que el fallo de un dispositivo o de su conexión
degrade la lectura de los cuatro. Con buses separados, la pérdida de un
sensor deja intactos a los restantes y resulta además localizable, pues
la línea afectada identifica al dispositivo.

La comunicación con la computadora de control se realiza por enlace
serial a \SI{115200}{baud}. El microcontrolador transmite
periódicamente una trama de texto delimitada por comas con cinco
campos: el promedio de las lecturas válidas, seguido de la temperatura
individual de cada uno de los cuatro sensores. Un sensor que no
responde no se omite de la trama ni conserva su último valor: su campo
se transmite como no numérico, de modo que la aplicación distingue la
pérdida de un sensor de una lectura legítima. El microcontrolador
acepta además un conjunto reducido de comandos que permiten verificar
la comunicación y suspender o reanudar la transmisión.

\subsection*{Consideraciones de firmware}

La adquisición de los cuatro sensores incorpora cuatro decisiones que
merecen mención.

La primera es el filtrado del valor de \SI{85}{\celsius}. Los
sensores DS18B20 devuelven ese valor como estado de reinicio tras el
encendido, antes de completar su primera conversión. Una lectura de
\SI{85}{\celsius} se descarta por corresponder a una conversión no
concluida y no a una medición válida.

La segunda es la detección de desconexión. Un sensor que deja de
responder produciría, con la implementación más directa, la repetición
indefinida de su última lectura válida, que es la peor forma de fallar:
el dato sigue pareciendo bueno. El firmware lleva por ello un contador
de fallos consecutivos por bus y, superado su umbral, emite un valor no
numérico en lugar de repetir el anterior. La recuperación es automática
en cuanto el sensor vuelve a responder.

La tercera es la identificación de los sensores. En un bus compartido,
la correspondencia entre el índice de lectura y el sensor físico
depende del orden de respuesta en el bus, de modo que una misma
posición del arreglo puede referirse a dispositivos distintos en
lecturas sucesivas. La topología de un sensor por bus elimina esa
ambigüedad sin necesidad de direccionamiento explícito: cada objeto de
lectura interroga una única línea, y el índice cero de esa línea
corresponde invariablemente al mismo dispositivo. La identificación
queda así resuelta por el cableado y no por el firmware.

La cuarta es el manejo del tiempo de conversión. A resolución de
\SI{12}{bits} la conversión de temperatura requiere aproximadamente
\SI{750}{\milli\second}. Bloquear el ciclo de ejecución durante ese
intervalo impide atender los comandos entrantes, por lo que la
conversión se solicita y se recoge de forma no bloqueante,
permitiendo que el firmware permanezca receptivo mientras la
conversión se encuentra en curso.

\subsection*{Colocación de los sensores}

Los cuatro sensores se disponen en las esquinas de la celda,
completamente sumergidos y próximos al fondo de la cavidad. Esta
disposición proporciona cobertura espacial del volumen de muestra y
permite detectar gradientes horizontales mediante la comparación
entre lecturas individuales.

Dos consideraciones acompañan esta elección.

La primera es de naturaleza acústica y resulta favorable. Con la
geometría establecida en la \cref{sec:celda}, una esquina se
encuentra a aproximadamente \SI{100}{\milli\meter} del eje del haz.
Una eventual reflexión acústica en el encapsulado metálico de un
sensor arribaría al hidrófono alrededor de \SI{68}{\micro\second}
después del disparo, muy por fuera de la ventana útil de
\SI{6}{\micro\second} establecida en la \cref{tab:tiempos_celda}. La
presencia de los sensores no perturba, por tanto, el intervalo de la
señal que se analiza.

La segunda es de naturaleza térmica y constituye una limitación. Con
un tirante de \SI{28.6}{\milli\meter} y el eje del haz situado a
aproximadamente \SI{20}{\milli\meter} del fondo, unos sensores
próximos al fondo miden el estrato inferior del líquido, que en
condiciones de reposo no coincide necesariamente con la temperatura
del volumen efectivamente iluminado. La magnitud de esta diferencia
depende del gradiente vertical presente y no fue cuantificada en este
trabajo. Se documenta aquí como fuente de incertidumbre sistemática y
se retoma en el \cref{cap:conclusiones}.

\subsection*{Calibración relativa entre sensores}

Los sensores DS18B20 presentan una dispersión de fábrica entre
unidades que, para el conjunto empleado, alcanza
\SI{0.32}{\celsius}. Dado que el criterio de disparo de la secuencia
automática descrito en el \cref{cap:resultados} opera sobre un umbral
de \SI{0.1}{\celsius} en torno a la temperatura objetivo, una
dispersión de esa magnitud excede el propio criterio de decisión y
debe corregirse. Se describe a continuación el procedimiento definido
para hacerlo, cuya aplicación quedó pendiente al cierre de este
trabajo.

La corrección adoptada es de carácter relativo y no absoluto. El
experimento no requiere que las lecturas coincidan con un patrón
trazable, sino que los cuatro sensores concuerden entre sí, puesto que
las condiciones de disparo se evalúan sobre transiciones de
temperatura y no sobre valores absolutos.

El procedimiento consiste en sumergir los cuatro sensores en un mismo
volumen de agua mantenido en agitación continua --- condición
indispensable, pues sin ella se mide la estratificación del líquido y
no la dispersión entre sensores --- y registrar sus lecturas en
varios puntos a lo largo del intervalo de operación. Se designa un
sensor como referencia y se calcula el desplazamiento medio de cada
uno de los tres restantes respecto de él. Las constantes de corrección
resultantes se incorporarían al firmware, de modo que la compensación
se aplicara antes del promediado.

Debe hacerse explícito que esta corrección no se encuentra implementada
en la versión del firmware empleada para las mediciones reportadas: los
cuatro sensores se promedian sin compensación entre unidades. En
consecuencia, el valor de temperatura registrado en cada captura
arrastra la dispersión de fábrica del conjunto, y el umbral de
\SI{0.1}{\celsius} que gobierna el criterio de disparo opera sobre
lecturas que no han sido igualadas entre sí. La consecuencia práctica
es acotada, pues dicho criterio se evalúa sobre el descenso monótono
del promedio y no sobre la concordancia entre sensores, pero conviene
tenerla presente al comparar lecturas individuales. La aplicación del
procedimiento se recoge en el \cref{cap:conclusiones} como tarea
pendiente.

% \begin{figure}[H]
%   \centering
%   \includegraphics[width=0.75\textwidth]{modulo_temperatura}
%   \caption{Módulo de medición de temperatura: microcontrolador
%            ESP32 WROOM-32 y sensores DS18B20.}
%   \label{fig:modulo_temperatura}
% \end{figure}

% \begin{figure}[H]
%   \centering
%   \includegraphics[width=0.75\textwidth]{diagrama_temperatura}
%   \caption{Diagrama de conexión de los cuatro buses 1-Wire
%            independientes y sus sensores.}
%   \label{fig:diagrama_temperatura}
% \end{figure}
